\documentclass[../../../include/open-logic-section]{subfiles}

\begin{document}

\olfileid{sth}{story}{rus}
\olsection{پارادوکس راسل (دوباره)}

در~\olref[sfr][][]{part} با نظریهٔ مجموعه‌های ساده‌انگارانه کار کردیم.
اما بنا بر تصوری \emph{بسیار} ساده‌انگارانه، مجموعه‌ها صرفاً مصداق‌های
محمول‌ها هستند. این اندیشهٔ ساده‌انگارانه اصل زیر را ایجاب می‌کند:

\begin{defish}
  \emph{اصل تصریح ساده‌انگارانه.} برای هر فرمول~$\phi$،
  $\Setabs{x}{\phi(x)}$ وجود دارد.
\end{defish}

این اصل، با همهٔ جذابیتش، به‌طور اثبات‌پذیر ناسازگار است. این امر را
در~\olref[sfr][set][rus]{sec} دیدیم، اما نتیجه چنان مهم و سرراست است
که ارزش تکرار لفظ‌به‌لفظ دارد.

\begin{thm}[پارادوکس راسل]
هیچ مجموعه‌ای~$R = \Setabs{x}{x \notin x}$ وجود ندارد.
\end{thm}

\begin{proof}
اگر~$R = \Setabs{x}{x \notin x}$ وجود داشته باشد، آنگاه
$R \in R$ اگر و تنها اگر~$R \notin R$، که تناقض است.
\end{proof}

راسل این نتیجه را در ژوئن ۱۹۰۱ کشف کرد. (اما پارادوکس را دقیقاً به
صورتی که اکنون ارائه کردیم بیان نکرد، زیرا نظریهٔ مجموعه‌های فرگه را،
چنان‌که در~\emph{Grundgesetze} آمده است، بررسی می‌کرد. در
\olref[blv]{sec} به این موضوع بازمی‌گردیم.) راسل در ۱۶ ژوئن ۱۹۰۲ به
فرگه نامه نوشت و ناسازگاری دستگاه او را توضیح داد. برای مکاتبات و اندکی
پیشینه، بنگرید به~\citet[pp.~124--8]{Heijenoort1967}.

شایان تأکید است که این اثبات دوخطی نتیجه‌ای از \emph{منطق محض} است.
البته در نمادگذاری~$\Setabs{x}{x \notin x}$ به‌طور ضمنی از یک اصل موضوع
(غیرمنطقی؟)، یعنی گسترش‌مندی، استفاده کردیم؛ زیرا
$\Setabs{x}{\phi(x)}$ بناست \emph{مجموعهٔ یگانهٔ} (بنا بر گسترش‌مندی)
اشیای~$\phi$ باشد، اگر چنین مجموعه‌ای وجود داشته باشد. اما می‌توانیم
حتی سایهٔ گسترش‌مندی را نیز با بیان نتیجه به صورت زیر کنار بگذاریم:
\emph{هیچ مجموعه‌ای وجود ندارد که اعضایش دقیقاً مجموعه‌های خودناعضو
باشند}. و این امر ارتباط ویژه‌ای با مجموعه‌ها ندارد. همان‌گونه که خود
راسل ملاحظه کرد، استدلالی دقیقاً مشابه شما را به این نتیجه می‌رساند:
\emph{هیچ مردی دقیقاً مردانی را اصلاح نمی‌کند که خود را اصلاح نمی‌کنند}.
یا: \emph{هیچ سگ پاگ دقیقاً سگ‌های پاگی را بو نمی‌کشد که خود را بو
نمی‌کشند}. و مانند آن. به‌طور شماتیک، صورت نتیجه صرفاً چنین است:
\[
\lnot \exists x \forall z(Rxz \liff \lnot Rzz).
\]
و این صرفاً یک قضیه (طرحواره) از منطق مرتبهٔ اول است. در نتیجه،
نمی‌توانیم فقط با دستکاری نظریهٔ مجموعه‌های خود از پارادوکس راسل
بگریزیم؛ این پارادوکس حتی پیش از آنکه \emph{به} نظریهٔ مجموعه‌ها برسیم
پدید می‌آید. اگر بناست از منطق مرتبهٔ اول (کلاسیک) استفاده کنیم، ناچاریم
\emph{بپذیریم} که مجموعه‌ای~$R = \Setabs{x}{x\notin x}$ وجود ندارد.

نتیجه این است: اگر بخواهید اصل تصریح ساده‌انگارانه را بپذیرید و درعین‌حال
از ناسازگاری \emph{بپرهیزید}، نمی‌توانید صرفاً \emph{نظریهٔ مجموعه‌ها}
را دستکاری کنید. در عوض باید \emph{منطق} خود را از نو بسازید.

البته نظریه‌های مجموعه‌ای با منطق‌های ناکلاسیک عرضه شده‌اند. اما آن‌ها،
دست‌کم، نامتعارف‌اند. رویکرد معیار به پارادوکس راسل آن است که آن را
اثباتی سرراست بر عدم وجود بدانیم و سپس بکوشیم بیاموزیم چگونه با آن
زندگی کنیم. ما همین رویکرد را دنبال خواهیم کرد.

\end{document}
